\documentclass[../main.tex]{subfiles}
\begin{document}

\chapter{Decision making under uncertainty}

\section{Markov decision problems}

\subsection{Basics}


A \emph{Markov decision process} is used to model the interaction of an agent and an environment.

It is described by a tuple $\left\langle \mathcal{S}, \mathcal{A}, p, p_{R}, p_{0} \right\rangle $, where $\mathcal{S}$ is a set of environment states
, $\mathcal{A}$ a set of actions the agent can take, $p$ is a \emph{transition model}, $p_{R}$ a \emph{reward model}, and $p_{0}$ the initial state distribution. The initial state $s_0 \sim  p_0$.

At time $t \ge  0$, the agent observes $s_{t} \in  \mathcal{S}$, follows a \emph{policy} $\pi$ to take an action $a_{t} \in \mathcal{A}$. The environment emits $r_{t} \in  \mathcal{R}$ and 
enters a new state $s_{t+1} \in  \mathcal{S}$.

We use $\pi(a|s)$ to denote the probability of choosing $a$ in state $s$. $\pi(s)$ is the conditional probability over $\mathcal{A}$ if the policy is stochastic.
The process at every step $t$ is called \emph{transition}.At time $t$, it consists of the tuple $(s_{t}, a_{t}, r_{t}, s_{t+1})$, where $a_{t} \sim  \pi(s_{t})$, $s_{t+1} \sim  p(s_{t}, a_{t})$ and $r_{t}\sim  p_{R}(s_{t}, a_{t}, s_{t+1})$.

Hence, the probability  of generating a trajectory $\tau$ of length $T$ can be written explicity as 
\[
p(\tau) = p_0(s_0) \prod_{t=0}^{T-1} \pi(a_{t}|s_{t}) p(s_{t+1} | s_t,a_{t}) p_{R}(s_{t}, a_{t}, s_{t+1}).
\]

Define the \emph{reward function} from $p_{R}$, 
\[
R(s,a) =  \mathbb{E}_{p(s'|s,a)} \left[ \mathbb{E} _{p_{R}(r|s,a,s')}[r] \right] .
\]

\paragraph{Partially observed MDPs}

The agent only sees a potentially noisy observation generated from $s_{t}$
\[
o_{t} \sim  p(\cdot |s_{t}, a_{t}).
\]
Now the policy is a mapping from data to actions, $a_{t} \sim  \pi(h_{t})$, where $h_{t} = (a_1,o_1, \ldots ,a_{t-1}, o_t)$


\paragraph{Episodes and returns}

The MDP describes how a trajectory $\tau$ is stochastically generated.

Let $\tau$ be a trajectory of length $T$, where $T$ may be $\infty$. We denote the return for the state at time $t$
\[
G_{t} = r_{t}  + \gamma r_{t+1} + \cdots + \gamma^{T-t-1} = \sum_{j=t}^{T-1}\gamma^{j-t}r_{j}.
\]
$G_{t}$ is called the \emph{reward-to-go}.

Obviously, $G_{t} = r_{t} + \gamma G_{t+1}$.


\paragraph{Value Function}

Let $\pi$ be a given policy. We define the \emph{state-value function}, or \emph{value function} as 
\[
V_{\pi}(s) = \mathbb{E}_{\pi}\left[ G_{0}|s_0=s \right]  = \mathbb{E}_{\pi} \left[ \sum_{t=0}^{\infty} \gamma^{t} r_{t} | s_{0} =s \right] .
\]
Similarly, we define \emph{action-value function}, or \emph{Q-function} as 
\[
Q_{\pi}(s,a) = \mathbb{E}_{\pi} \left[ G_0|s_0=s, a_0=a \right]  = \mathbb{E}_{\pi} \left[ \sum_{t=0}^{\infty}\gamma^{t} r_{t} | s_0=s, a_0=a \right] .
\]
We have $V_{\pi}(s) = \mathbb{E}_{\pi(a|s)} \left[ Q_{\pi}(s,a) \right] $.
Define the \emph{advantage function} as 
\[
A_{\pi}(s,a) = Q_{\pi}(s,a) - V_{\pi}(s).
\]
This shows the benefit of picking action $a$ in state $s$ then switching to policy $\pi$.


\paragraph{Optimal value functions and policies}

Optimal policy $\pi_{*}$, optimal value function $V_{*}$.

\emph{Bellman's optimality equations}:
\begin{equation}
\begin{aligned}
V_{*}(s) &= \max_{a} R(s,a) + \gamma \mathbb{E}_{p(s'|s,a)} \left[ V_{*}(s') \right]  \\
Q_{*}(s,a) &= R(s,a) + \gamma \mathbb{E}_{p(s'|s,a)} \left[ \max_{a'} Q_{*}(s',a')\right] 
\end{aligned}
\end{equation}

Given an optimal value function, we can derive an optimal policy using
\[
\begin{aligned}
\pi_{*}(s) &= \mathop{\arg\max}\limits_{a}Q_{*}(s,a) \\
&= \mathop{\arg\max}\limits_{a} \left[ R(s,a)  + \gamma \mathbb{E}_{p(s'|s,a)}\left[ V_{*}(s') \right] \right] .
\end{aligned}
\]

\subsection{Planning in MDPs}

How to compute an optimal policy when the MDP is known.


\paragraph{Value iteration}












\end{document}
